\documentclass[10pt,a4paper]{amsart}
\usepackage[margin=19mm]{geometry}
\usepackage{amsmath,amssymb,mathtools}
\usepackage[scr=rsfs,cal=euler]{mathalfa}
\usepackage{xfrac}
\usepackage[unicode,hidelinks,bookmarks=false]{hyperref}
\newcommand{\ie}{i.e.\ }
\newcommand{\cf}{cf.\ }
\newcommand{\resp}{respectively\ }
\newcommand{\Subsection}{\subsection}
\providecommand{\nobreakdash}{\nobreak\hbox{-}}
\setlength{\emergencystretch}{2em}
\multlinegap0pt
\usepackage{amssymb}
\usepackage{enumerate}
\newtheorem{theorem}{Theorem}
\title{On transcendence of non-periodic continued fractions associated with modular forms and arithmetic functions}
\author{Tapas Chatterjee, Sagar Mandal}
\date{Source: arXiv:2602.13300v1 (CC BY 4.0)}
\begin{document}
\maketitle

\noindent\emph{Source and licence.} On transcendence of non-periodic continued fractions associated with modular forms and arithmetic functions. Version arXiv:2602.13300v1 (CC BY 4.0), distributed by arXiv under the Creative Commons Attribution 4.0 International licence, http://creativecommons.org/licenses/by/4.0/, as recorded in the arXiv licence metadata for this exact version and shown on the abstract page https://arxiv.org/abs/2602.13300v1. Full licence notice: \texttt{licenses/arxiv-2602.13300-CC-BY-4.0.txt}.\par\smallskip

\noindent\textit{Editorial scope.} Counted unit: the article's theorem thm1.8 (source0.tex lines 158--160). The article's complete original proof of it is delivered with it (source0.tex lines 223--249, one \texttt{proof} environment of 27 lines, which the article places away from the statement). No mathematics is written, repaired or re-derived: the statement and the proof are the article's own lines, in the article's own notation, and no retained line is substituted or re-flowed.  No further source-local context is needed: every cross-reference of every retained line resolves inside the two delivered spans.  Nothing was omitted from any retained span, so the package carries no omission record.  No retained line contains a citation, so the wrapper carries no bibliography and no bibliography entry.  No retained line contains a figure environment, an image-inclusion command or drawing code, so no image is delivered and the package declares no asset dependency.  Editorial formatting only, in addition: an AMS \texttt{amsart} preamble that restates the article's own declarations and macros used by the retained lines, namely the environments \texttt{theorem} on separate counters, together with the standard packages those lines need.  The retained environments print in this document as Theorem 1 in order of appearance, whereas the article prints the same items with the numbering of its own full document.  The complete native source of the article as distributed in the arXiv e-print for this exact version is bundled unchanged as \texttt{sources/194628/source0.tex}.
\begin{theorem}\label{thm1.8}
 Let $f$ be an arithmetic function and $m$ be a positive integer such that $m\nmid f(p)$ for infinitely many primes $p$ and $m\mid f(p)$ for primes $p\equiv 1 \pmod{K(m)}$, where $K(m)$ is some positive integer-valued function of $m$. If $f$  has the property that $f(p)\mid f(n)$ whenever $p\mid n$, where $p$ is a prime, then $\{f(n) \pmod{m}\}_{n\geq1}$ is non-periodic.   
\end{theorem}

\begin{proof}[\textbf{Proof of Theorem \ref{thm1.8}}]
Assume, for the sake of contradiction, that the sequence $\{f(n) \pmod{m}\}_{n\geq 1}$ is periodic. Then there exist fixed positive integers $N,L$ such that
\begin{align}\label{A}
f(n+L)\equiv f(n) \pmod{m}\quad \text{for all~}n\geq N.    
\end{align}

In particular, choose $n=p_N\geq N$, the smallest prime such that $f(p_N)\not\equiv 0\pmod{m}$  (existence is guaranteed by the fact that $m\nmid f(p)$ for infinitely many primes $p$), then from (\ref{A}) we get
\begin{align}\label{1}
f(p_N+jL)\equiv f(p_N)\not\equiv 0\pmod{m}\quad \text{for all~}j\in\mathbb{N}.    
\end{align}
Setting $j=p_NK(m)j_N$ in (\ref{1}), where $j_N$ is some positive integer, we obtain
\begin{align}\label{2}
f(p_N(1+K(m)j_NL))\equiv f(p_N)\not\equiv 0\pmod{m}.    
\end{align}
Since $\gcd(1,p_NK(m)L)=1$, by the Dirichlet prime number theorem, the following sequence
$$1, \quad 1+p_NK(m)L, \quad 1+2p_NK(m)L,\quad \ldots,\quad 1+j_N'p_NK(m)L,\quad \ldots$$
contains infinitely many primes. Without loss of generality, assume that $1+j_N'p_NK(m)L$ is the smallest prime in this sequence. Then choosing $j_N=j_N'p_N$ in (\ref{2}) we get
\begin{align}\label{eq3}
f(p_N(1+K(m)j_N'p_NL))\equiv f(p_N)\not\equiv 0\pmod{m}.    
\end{align}
Since $1+K(m)j_N'p_NL$ is a prime and satisfies $1+K(m)j_N'p_NL\equiv 1 \pmod{K(m)}$, we obtain  $m\mid f(1+K(m)j_N'p_NL)$. Moreover, using the fact that $f(p)\mid f(n)$ whenever $p\mid n$ and observing that $1+K(m)j_N'p_NL\mid p_N(1+K(m)j_N'p_NL$ we deduce  $f(1+K(m)j_N'p_NL)\mid f(p_N(1+K(m)j_N'p_NL))$. Consequently, it follows that
\begin{align}\label{1*}
 0\equiv f(p_N(1+K(m)j_N'p_NL)\equiv f(p_N)\not\equiv 0\pmod{m},   
\end{align}

which is a contradiction. This completes the proof.
\end{proof}

\end{document}
